\documentclass[11pt]{article}
\usepackage[margin=1in]{geometry}
\usepackage{amsmath,amssymb,amsthm,mathtools}
\IfFileExists{lmodern.sty}{\usepackage{lmodern}}{}
\usepackage[expansion=false]{microtype}
\usepackage{booktabs,array}
\usepackage{enumitem}
\usepackage{xcolor}
\usepackage[colorlinks=true,linkcolor=blue!50!black,citecolor=green!35!black,urlcolor=blue!50!black]{hyperref}
\usepackage[nameinlink]{cleveref}
\newcolumntype{R}[1]{>{\raggedright\arraybackslash}p{#1}}

\newtheorem{theorem}{Theorem}[section]
\newtheorem{proposition}[theorem]{Proposition}
\newtheorem{lemma}[theorem]{Lemma}
\newtheorem{corollary}[theorem]{Corollary}
\newtheorem{conjecture}[theorem]{Conjecture}
\theoremstyle{definition}
\newtheorem{definition}[theorem]{Definition}
\newtheorem{example}[theorem]{Example}
\newtheorem{question}[theorem]{Question}
\theoremstyle{remark}
\newtheorem{remark}[theorem]{Remark}

\newcommand{\D}{\mathsf{D}}
\newcommand{\Kl}{\mathrm{Kl}}
\newcommand{\Cpl}{\operatorname{Cpl}}
\newcommand{\TV}{\mathrm{TV}}
\newcommand{\Hyp}{\mathrm{Hyp}}
\newcommand{\Sh}{\mathrm{Sh}}
\newcommand{\C}{\mathsf{C}}
\newcommand{\supp}{\operatorname{supp}}
\newcommand{\id}{\mathrm{id}}

\title{Conditionals, Couplings and Toric Geometry:\\
Kernels, the Coupling Nerve, Cumulants and a Fukaya Programme\\[4pt]
\large Revised draft 3}
\author{Vinay K. Awasthi\\
Department of Applied Mathematics and Statistics\\
Johns Hopkins University\thanks{Verification scripts are listed in \cref{sec:code}; the worked examples are
in the companion \cite{Awasthi-Family}.}}
\date{October 6, 2026}

\begin{document}
\maketitle

\begin{abstract}
We organize the conditional and coupling structure of the finite distribution monad $\D$ into four layers
and separate what is proved from what is conjectural.
\begin{enumerate}[label=(\arabic*),leftmargin=*]
\item \emph{The base.} Markov kernels are the 1-cells of the Kleisli category, and couplings are 2-cells whose
  vertical composition is the gluing lemma. The natural whiskerings are not functorial, so the couplings do not
  form a 2-category. Their image under the cost map does: a strict, graded, Lawvere-enriched 2-category.
\item \emph{The fibre.} The base-relative posterior tower is a functor on the Kleisli category. Its category of
  elements is a discrete opfibration over it, and every grade of the shadow bounds the distance between
  transported contexts at every level of the tower.
\item \emph{The coupling nerve.} Couplings of all arities form the simplicial object $\D(Y^{\bullet+1})$, with
  marginalization as faces. Every 2-horn fills. The fillers form a product of coupling polytopes, and the
  maximum-entropy filler is the one whose third face is the vertical composite. For $n\ge3$ the faces of an
  $n$-horn form a cyclic cover, so filling can fail. The failure is measured by the contextual fraction, and
  $\D(Y^{\bullet+1})$ is not a Kan complex, although $Y^{\bullet+1}$ is.
\item \emph{Linearization.} Centred couplings are second mixed cumulants, and whiskering defects are the
  Boolean cumulants of whiskering viewed as a noncommutative expectation. Into the square-zero cone of the
  circulation ideal, whiskering is an $A_\infty$-functor with exactly two components. Its cohomology keeps
  only the marginals, and the coupling data are homotopical.
\end{enumerate}
On top of these layers sits the geometry. Coupling polytopes of nondegenerate marginals are Delzant, so they
are moment polytopes of toric manifolds whose Kähler class is the product coupling and whose canonical metric
restricts to the Fisher metric. Their walls and chambers, monotone cases, Floer theory and Bridgeland
stability conditions are treated as before. Conjectures are stated with correct types, and a companion paper
\cite{Awasthi-Family} computes every construction explicitly.
\end{abstract}

\tableofcontents

%=====================================================================
\section{Introduction}\label{sec:intro}
%=====================================================================

Two structures sit on top of the finite distribution monad $\D$. The base-relative posterior tower
\cite{Awasthi-Tower} fixes an outcome base $B$ and solves a \emph{strict integrability} problem: it carries a
mixed distributive law satisfying all four Beck axioms. The conditional/coupling structure
\cite{Awasthi-Attribution} lets the base vary along Markov kernels and compares kernels by couplings: it solves
a \emph{transport} problem. This paper assembles them, together with a simplicial and a homotopical layer,
into one framework.

\begin{center}\small
\begin{tabular}{@{}R{0.17\textwidth}R{0.42\textwidth}R{0.33\textwidth}@{}}
\toprule
Layer & Objects and structure & Status\\
\midrule
I.\ Base & finite sets, kernels, coupling 2-cells; the graded shadow & strict 2-category on grades only
(\cref{sec:base})\\
II.\ Fibre & the tower over each base; transport along kernels & discrete opfibration; grades act on fibres
(\cref{sec:fibre})\\
III.\ Nerve & $\D(Y^{\bullet+1})$; horns, fillers, boundary fibres & 2-horns fill; higher horns obstructed by
contextuality (\cref{sec:nerve})\\
IV.\ Linearization & circulations, cumulants, cones & whiskering is a two-component $A_\infty$-functor into a cone
(\cref{sec:linear})\\
Geometry & coupling polytopes as moment polytopes; stability & proved in the nondegenerate case
(\cref{sec:toric,sec:imported,sec:stability})\\
Programme & families over chambers; Fukaya-type functor & conjectural (\cref{sec:conjectural})\\
\bottomrule
\end{tabular}
\end{center}

The layers interlock. A 2-cell is a 1-simplex of the nerve, and vertical composition is the canonical filler of
an inner 2-horn. The tangent spaces of the filler polytopes are circulations, and those are where the cumulants
of layer IV live. The tower of layer II is transported along the 1-cells of layer I, and the grades of layer I
bound the transport. The polytopes of layer III are, in the nondegenerate case, moment polytopes, and that is
where the geometry enters. What does \emph{not} hold is equally important. The couplings themselves do not form a
2-category, the nerve is not Kan, and the toric picture stops at boundary fibres of 3-way tables.

%=====================================================================
\section{Layer I: kernels, couplings and the graded shadow}\label{sec:base}
%=====================================================================

$\D$ is the finitely supported distribution monad on $\mathbf{Set}$ and $\Kl(\D)$ its Kleisli category.
A 1-cell $k\colon B\rightsquigarrow B'$ is a Markov kernel $k\colon B\to\D(B')$, acting on states by
$k^\#(p)(b')=\sum_bp(b)k(b)(b')$. The results of this section are proved in
\cite[\S7]{Awasthi-Attribution}; we restate them for reference.

\begin{proposition}[1-cells {\cite[Props.~7.2, 7.3]{Awasthi-Attribution}}]\label{prop:1cells}
$k^\#$ is affine, $(h\circ k)^\#=h^\#k^\#$, $k^\#\circ\mu=\mu\circ\D(k^\#)$, and
$\TV(k^\#p,k^\#q)\le\TV(p,q)$. Attribution sets are transported along $k$, and precision can only be lost.
\end{proposition}

\begin{definition}[Coupling 2-cells]\label{def:coupling}
For 1-cells $f,g\colon X\rightsquigarrow Y$, a coupling 2-cell $\alpha\colon f\Rightarrow g$ is a kernel
$\alpha\colon X\to\D(Y\times Y)$ with marginals $f(x)$ and $g(x)$. Its cost is
$c(\alpha)=\max_x\Pr_{\alpha(x)}(y\ne y')$; the identity is the diagonal coupling; the vertical
composite glues along the common marginal,
$(\beta\cdot\alpha)(x)(y,y'')=\sum_{y'}\alpha(x)(y,y')\beta(x)(y',y'')/g(x)(y')$.
\end{definition}

\begin{theorem}[Witnesses and shadow {\cite[Thms.~7.6, 7.8]{Awasthi-Attribution}}]\label{thm:witness}
\leavevmode
\begin{enumerate}[label=(\roman*)]
\item For fixed $X,Y$, 1-cells and coupling 2-cells form a category under gluing, strictly
  associative and unital. Cost is subadditive and $\min_\alpha c(\alpha)=d(f,g)=\max_x\TV(f(x),g(x))$.
\item Whiskering by a deterministic precomposition is strictly functorial. Whiskering by a stochastic
  precomposition (mixing couplings) or by a postcomposition (copying on equal pairs, independent on
  unequal pairs) preserves identities, marginals and the cost bound, but is not functorial.
\item Graded 2-cells $f\Rightarrow_\varepsilon g$, $\varepsilon\ge d(f,g)$, with grades adding, form a strict
  2-category on kernel 1-cells, and $d$ makes $\Kl(\D)$ a Lawvere-enriched category.
\end{enumerate}
\end{theorem}

\begin{remark}[Witness layer and shadow]\label{rem:shadow}
The cost map $\alpha\mapsto c(\alpha)$ sends the witness layer onto the graded shadow. It preserves identities and
sends composites to grades no larger than the sum. The shadow is an image of the witness layer, not a
substructure of it, and it is the only strict 2-category in sight: with couplings as 2-cells, (ii) rules out a
2-category or a bicategory with these whiskerings. The polytopes $\Cpl(f(x),g(x))$ are the spaces of 2-cells, not
further cells; their geometry is the subject of \cref{sec:nerve,sec:toric}.
\end{remark}

%=====================================================================
\section{Layer II: the tower as a fibration}\label{sec:fibre}
%=====================================================================

The base-relative tower over $B$ has context algebra $A_B=\D B\times\D^2B\times\D^3B\times\Sh(B)$,
reader comonad $\C^B(Y)=Y\times A_B$ and strict law $\lambda^B$ \cite[\S9]{Awasthi-Tower}. With
$\Hyp=\Theta\times\D(-)$, a kernel $k\colon B\to\D(B')$ acts on hypotheses by
$\Hyp(k)(\theta,\rho)=(\theta,k^\#\rho)$, and we put
$A_k=k^\#\times\D(k^\#)\times\D^2(k^\#)\times\D(\Hyp k)$.

\begin{theorem}[The tower is a functor on the Kleisli category {\cite[Prop.~9.14, Cor.~9.15, Prop.~9.16]{Awasthi-Tower}}]\label{thm:fibre}
\leavevmode
\begin{enumerate}[label=(\roman*)]
\item $A_k$ is a morphism of $\D$-algebras, preserves coherent contexts, commutes with the canonical
  section $\sigma$, and $A_{h\circ k}=A_hA_k$; for deterministic $k$ it is the earlier base change.
\item $\id\times A_k$ is a comonad morphism $\C^B\Rightarrow\C^{B'}$ commuting with the laws $\lambda^B$,
  $\lambda^{B'}$.
\item Admissibility data, coherent lifts and the effective sets are transported; certificates pull
  back along $f'\mapsto k\rhd f'$; failure levels can only rise and the secondary obstruction can only
  fall.
\item Conditioning on an outcome observed after the channel equals conditioning on the pulled-back
  predicate $k\rhd1_{e'}$, then changing base.
\end{enumerate}
\end{theorem}

\begin{definition}[The total category]\label{def:total}
Let $\mathcal E$ be the category whose objects are pairs $(B,z)$ with $B$ a finite set and $z\in A_B^{\mathrm{coh}}$ a
coherent context, and whose morphisms $(B,z)\to(B',z')$ are kernels $k$ with $A_k(z)=z'$. Let $\pi\colon\mathcal E\to\Kl(\D)$
forget $z$.
\end{definition}

\begin{proposition}[Split opfibration]\label{prop:opfib}
$\pi$ is a discrete opfibration, and in particular a split opfibration. The cocartesian lift of $k\colon B\to B'$ at
$(B,z)$ is $k\colon(B,z)\to(B',A_kz)$. Through $\id\times A_k$, the comonads $\C^B$ and the laws $\lambda^B$ are carried
along the lifts.
\end{proposition}
\begin{proof}
$\mathcal E$ is the category of elements of the functor $B\mapsto A^{\mathrm{coh}}_B$, $k\mapsto A_k$, which is a functor by
\cref{thm:fibre}(i); categories of elements of set-valued functors are discrete opfibrations. The statement about
comonads and laws is \cref{thm:fibre}(ii).
\end{proof}

\begin{proposition}[Grades act on fibres]\label{prop:gradefibre}
Let $f,g\colon B\to\D(B')$ and let $z=(\rho,\Pi,\Omega,\Sigma)$ be a coherent context. Put $\mathrm{TV}$ on $\D B'$, the
Kantorovich metric $W_1$ over $\mathrm{TV}$ on $\D^2B'$, the Kantorovich metric over that on $\D^3B'$, and on hypotheses the
metric that is $\TV$ of predictives for equal descriptors and $1$ otherwise. Then every component of $A_f(z)$ and
$A_g(z)$ is within $d(f,g)$. So a graded 2-cell $f\Rightarrow_\varepsilon g$ gives $\varepsilon$-close transported contexts at
every level.
\end{proposition}
\begin{proof}
$\TV(f^\#\rho,g^\#\rho)\le\sum_b\rho(b)\TV(f(b),g(b))\le d(f,g)$ by convexity. For $\Pi$, push $\Pi$ forward along
$\rho\mapsto(f^\#\rho,g^\#\rho)$. This is a coupling of $\D(f^\#)\Pi$ and $\D(g^\#)\Pi$ of cost $\mathbb E_\Pi\TV(f^\#\rho,g^\#\rho)\le d(f,g)$.
Level three and the shadow are the same argument one level up. The companion verifies the bound in a worked
example, where all three levels give $\tfrac2{15}\le\tfrac16$ \cite[\S11]{Awasthi-Family}.
\end{proof}

The proof of \cref{thm:fibre} uses one fact beyond naturality: Kleisli extension is a morphism of free
$\D$-algebras. All other components of $A_k$ are pushforwards along functions. So the strict integrability of
each fibre and the transport between fibres are compatible.

%=====================================================================
\section{Layer III: the coupling nerve}\label{sec:nerve}
%=====================================================================

\begin{definition}[The coupling nerve]\label{def:nerve}
For a finite set $Y$, let $N_k(Y)=\D(Y^{k+1})$, with face maps $d_i$ that marginalize out the $i$th coordinate and
degeneracies $s_i$ that push forward along the diagonal repeating it. Since $Y^{\bullet+1}$ is a simplicial set (the
Čech nerve of $Y\to1$) and $\D$ is a functor, $N(Y)=\D(Y^{\bullet+1})$ is a simplicial set. The relative version over a
set $X$ has $k$-simplices the kernels $X\to\D(Y_0\times\dots\times Y_k)$, with all structure taken pointwise in $x$.
\end{definition}

A 1-simplex of $N(Y)$ is a coupling of its two vertices. A coupling 2-cell $\alpha\colon f\Rightarrow g$ between kernels
$X\rightsquigarrow Y$ is an $X$-indexed 1-simplex with vertices $f$ and $g$. A horn $\Lambda^n_k$ in $N$ is a compatible family of
laws on the $n$ faces other than the $k$th, and a filler is an $n$-simplex with those faces. In probabilistic terms
a horn is an empirical model on the cover of $\{0,\dots,n\}$ by those faces, and a filler is a global section
\cite{AbramskyBrandenburger2011}.

\begin{theorem}[2-horns fill, and composition is the canonical filler]\label{thm:2horn}
Let $p_{01}$ and $p_{12}$ be compatible laws on $Y_0\times Y_1$ and $Y_1\times Y_2$ with common marginal $p_1$, the data of
an inner horn $\Lambda^2_1$.
\begin{enumerate}[label=(\roman*)]
\item The fillers form the product, over $x_1$ with $p_1(x_1)>0$, of the transportation polytopes of the slices
  $p(\cdot,x_1,\cdot)$ with marginals $p_{01}(\cdot,x_1)$ and $p_{12}(x_1,\cdot)$. If every slice has full support and
  nondegenerate marginals, the filler polytope is Delzant.
\item The gluing filler $p^{\mathrm{gl}}=p_{01}\,p_{12}/p_1$ is the unique filler of maximal entropy, and its face $d_1$ is the
  vertical composite of the two couplings.
\item Any filler differs from $p^{\mathrm{gl}}$, slice by slice, by $p_1(x_1)\operatorname{Cov}(1_{\{x_0\}},1_{\{x_2\}}\mid x_1)$, a
  circulation of the slice.
\end{enumerate}
The outer horns $\Lambda^2_0$ and $\Lambda^2_2$ fill in the same way, by gluing over the shared vertex.
\end{theorem}
\begin{proof}
(i) Each marginal constraint involves one value of $x_1$, so the constraints decouple into slices. Each slice is a
two-way transportation polytope, Delzant under the stated conditions (\cref{thm:delzant}), and a product of
Delzant polytopes is Delzant with respect to the product lattice. (ii) $\log p^{\mathrm{gl}}$ is a sum of a function of
$(x_0,x_1)$ and a function of $(x_1,x_2)$, so it is orthogonal to every tangent vector of the filler polytope. So
$p^{\mathrm{gl}}$ is a critical point of the strictly concave entropy on the polytope, hence its unique maximizer.
Equivalently, $H(X_0X_1X_2)\le H(X_0X_1)+H(X_1X_2)-H(X_1)$, with equality exactly under conditional independence.
Summing $p^{\mathrm{gl}}$ over $x_1$ gives the gluing formula of \cref{def:coupling}. (iii) is a direct computation.
\end{proof}

So the 2-cells of layer I are the 1-simplices of the nerve, and their vertical composition is the canonical
filling of inner 2-horns. The other fillers are conditional cumulants, which is the link to layer IV.

\begin{theorem}[Higher horns]\label{thm:higherhorn}
Let $n\ge3$.
\begin{enumerate}[label=(\roman*)]
\item The $n$ faces of $\Lambda^n_k$, viewed as subsets of $\{0,\dots,n\}$, form a cover that is not acyclic. For $n=2$ it is
  acyclic.
\item Hence, for $\lvert Y\rvert\ge2$, there are compatible horn data in $N(Y)$ with no filler.
\item For compatible horn data, the minimal unmodeled weight with the deterministic global assignments as
  sources, each face weighted $1/n$, equals the contextual fraction of the horn, and it vanishes exactly when
  the horn fills. For the parity horns on two outcomes (uniform on assignments of fixed parity on each face, with
  exactly one face odd) it is $\tfrac12$ for every 3-horn and $\tfrac34$ for every 4-horn.
\end{enumerate}
\end{theorem}
\begin{proof}
(i) Use Graham's reduction for acyclic hypergraphs \cite{BFMY1983}: repeatedly delete vertices in a single edge and
edges contained in others; the hypergraph is acyclic if and only if this empties it. In $\Lambda^n_k$ the vertex $k$
lies in all $n$ faces and every other vertex lies in $n-1$ of them. No face contains another, so for $n\ge3$ no
reduction step applies. For $n=2$ the two faces share one vertex and reduce away. (ii) By Vorob'ev's theorem
\cite{Vorobev1962}, every compatible family extends to a global law exactly on acyclic covers. Data on two outcomes
embed into any $Y$ with $\lvert Y\rvert\ge2$. (iii) A filler is a global section, so this is
\cite[Thm.~7.12]{Awasthi-Attribution} applied to the horn's cover. The parity values are computed with exactly verified
certificates in \cite[\S10]{Awasthi-Family}.
\end{proof}

\begin{corollary}[The nerve is not Kan]\label{cor:notkan}
For $\lvert Y\rvert\ge2$, $N(Y)=\D(Y^{\bullet+1})$ fills every 2-horn but not every $n$-horn for $n\ge3$. In particular it is not
a Kan complex, although $Y^{\bullet+1}$ is. The monad $\D$ does not preserve horn filling, because it does not turn
compatible families over cyclic covers into global laws, and the obstruction is the contextual fraction.
\end{corollary}

\begin{proposition}[Boundary fibres]\label{prop:boundaryfibre}
Prescribing all three faces of a 2-simplex, that is all three 2-margins of a 3-way table on $Y_0\times Y_1\times Y_2$, gives
a planar 3-way transportation polytope of dimension $(n_0-1)(n_1-1)(n_2-1)$ when it is nonempty and generic. It can be empty
even for compatible faces, because the cover by the three edges of a triangle is cyclic.
\begin{enumerate}[label=(\roman*)]
\item For $n_0=2$ it is a capacitated two-way transportation polytope, whose constraint matrix is totally unimodular,
  so it is Delzant whenever it is simple.
\item Every rational polytope is a planar $r\times s\times3$ transportation polytope \cite{DeLoeraOnn2004}, so boundary fibres are
  not toric in general. Already for $3\times3\times3$, generic instances are simple but not smooth: they are moment polytopes
  of toric orbifolds \cite{LermanTolman1997}, not of manifolds \cite[\S10]{Awasthi-Family}.
\end{enumerate}
\end{proposition}
\begin{proof}
For $n_0=2$ the second layer is $m-y$, where $y$ is the first layer, $m$ is the $(1,2)$-margin, and $y$ has fixed row and column
sums with $0\le y\le m$. A bipartite incidence matrix stacked with an identity is totally unimodular, so edges are $\pm1$
cycle vectors and the argument of \cref{thm:delzant} applies at simple vertices. (ii) combines the cited theorem with
the computation in the companion.
\end{proof}

So the toric regime of the nerve is exactly the 1-simplices and the 2-horns. Boundary fillings leave it, and higher
horns can be empty.

%=====================================================================
\section{Layer IV: cumulants and the square-zero cone}\label{sec:linear}
%=====================================================================

\begin{proposition}[Couplings are second cumulants]\label{prop:kappa2}
For $\alpha\in\Cpl(f,g)$, $\alpha-f\otimes g$ is the cross-covariance $\kappa_2(1_{\{y=i\}},1_{\{y'=j\}})$ of the indicator vectors, a
circulation. The product coupling is the point where all mixed cumulants vanish, and it is the only point of
rank one, the Birch point of the independence model \cite{PachterSturmfels2005}. For level-two couplings of
second-order states,
\[
\mu^*(\Pi)=\mu_BP\otimes\mu_BQ+\operatorname{Cov}_\Pi(\rho,\rho'),
\]
so flattening keeps only the second mixed cumulant of the two random distributions.
\end{proposition}

Fix $f\in\D(Y)$ and let $A_f$ be the matrices on $Y\times Y$ whose row and column sums are multiples of $f$. With the
gluing product $a\cdot b=a\operatorname{diag}(f)^{-1}b$ and unit $\operatorname{diag}(f)$, $A_f$ is an associative unital algebra, spanned by
the couplings of $f$ with itself. The circulations $I_f$ form a two-sided ideal, and $A_f/I_f\cong\mathbb R$ records the total
mass. Whiskering by a kernel $h$ is a linear unital map $h_*\colon A_f\to A_{h^\#f}$, which need not be multiplicative: a
noncommutative expectation in the sense of homotopy probability theory \cite{HPT1,HPT2}. Its Boolean cumulants
\cite{SpeicherWoroudi1997} are $b_1=h_*$, $b_2(a,b)=h_*(a\cdot b)-h_*a\cdot h_*b$, and so on over interval partitions. The
whiskering defect of \cref{thm:witness}(ii) is $-b_2$.

\begin{theorem}[Whiskering is a two-component $A_\infty$-functor]\label{thm:sqzero}
Let $C=A_{h^\#f}\oplus\varepsilon I_{h^\#f}$, with $\varepsilon$ of degree $-1$, $\varepsilon^2=0$ and $d(\varepsilon x)=x$. Then $C$ is a differential graded
algebra with $H^0(C)=A_{h^\#f}/I_{h^\#f}\cong\mathbb R$ and $H^{-1}(C)=0$: it is the mapping cone of $I\hookrightarrow A$. The maps
\[
F_1=h_*,\qquad F_2=\varepsilon\,b_2,\qquad F_n=0\ (n\ge3)
\]
form an $A_\infty$-functor $A_f\to C$.
\end{theorem}
\begin{proof}
Every $b_2$ is a circulation, because both composites couple the same marginals, so $F_2$ lands in $\varepsilon I$. $C$ is the
trivial extension of $A$ by the bimodule $\varepsilon I$, and $d$ is a derivation because $I$ is an ideal. Arity two of the functor
equation is $dF_2=b_2$, true by definition. In arity three, with $F_3=0$, it reads
$b_2(ab,c)-b_2(a,bc)=h_*a\cdot b_2(b,c)-b_2(a,b)\cdot h_*c$. Both sides equal $h_*a\cdot h_*(bc)-h_*(ab)\cdot h_*c$, which says that $b_2$
is the Hochschild coboundary of $h_*$. In arity $n\ge4$ every term involves some $F_k$ with $k\ge3$ or a product
$F_2\cdot F_2\in\varepsilon^2I=0$.
\end{proof}

\begin{remark}
This settles, in the minimal model, the question of the previous draft whether whiskering is an $A_\infty$-functor.
The higher Boolean cumulants are nonzero \cite[\S9]{Awasthi-Family}, but they are not needed here. They appear as
higher components only in larger models such as the free cone, as in the bar-construction description of
cumulants. The defects span all of $I$ for three and four points \cite[\S10]{Awasthi-Family}, so $H^0$ of the cone
generated by the defects alone is also $A/I$.
\end{remark}

\paragraph{The four layers together.} Under the cone, all coupling information is homotopical: $H^0$ keeps only
masses, so the 1-cells of layer I, while the couplings live in degree $-1$, in the circulations. These are the
tangent spaces of the polytopes of layer III. The cumulants of \cref{prop:kappa2}, the defects of \cref{thm:sqzero}
and the deviations of 2-horn fillers from the canonical filler (\cref{thm:2horn}(iii)) are all circulations, and
the Fisher metric of \cref{prop:guillemin} measures them. Layer II carries the tower along the same kernels, with
the grades of layer I as bounds.

%=====================================================================
\section{Coupling polytopes are moment polytopes}\label{sec:toric}
%=====================================================================

\subsection{The polytope and its lattice}

Let $f\in\D(Y)$ and $g\in\D(Y')$ have full support, with $\lvert Y\rvert=n$, $\lvert Y'\rvert=m$. (For
coupling 2-cells $Y=Y'$; rectangular polytopes arise when gluing contexts
\cite[Prop.~7.11]{Awasthi-Attribution}.) The \emph{coupling polytope} is
\[
\Cpl(f,g)=\Big\{\alpha\in\mathbb{R}_{\ge0}^{Y\times Y'}:\ \textstyle\sum_{y'}\alpha(y,y')=f(y),\
\sum_y\alpha(y,y')=g(y')\Big\}.
\]
Its tangent space is the space of \emph{circulations}, matrices with zero row and column sums. The
integer circulations form a lattice $\Lambda$ of rank $(n-1)(m-1)$; a basis is
$E_{ij}-E_{i,m}-E_{n,j}+E_{n,m}$ for $i<n$, $j<m$.

\begin{proposition}[Interior and Fisher metric]\label{prop:interior}
$\Cpl(f,g)$ has dimension $(n-1)(m-1)$, its relative interior is the set of strictly positive
couplings, and on it the Fisher metric $\mathcal{I}_\alpha(v,w)=\sum v(y,y')w(y,y')/\alpha(y,y')$ is
positive definite on circulations.
\end{proposition}
\begin{proof}
The product coupling $f\otimes g$ is strictly positive, so no inequality $\alpha\ge0$ is an implicit
equality, and the relative interior of the polytope in its affine hull is $\{\alpha>0\}$. There
$\mathcal{I}_\alpha(v,v)=\sum v^2/\alpha>0$ for $v\ne0$.
\end{proof}

\begin{remark}
The first draft argued nondegeneracy through a radical and then distinguished ``the Fisher metric''
from ``the Fisher--Rao metric''. These are the same metric, the Fisher information metric. The
square-root embedding is a way of computing its geodesic distance \cite[Prop.~8.1]{Awasthi-Attribution}.
\end{remark}

\subsection{Nondegeneracy implies Delzant}

\begin{definition}[Walls]\label{def:walls}
The pair $(f,g)$ is \emph{degenerate} if $\sum_{y\in I}f(y)=\sum_{y'\in J}g(y')$ for some
nonempty proper subsets $I\subsetneq Y$ and $J\subsetneq Y'$. These equations define hyperplanes,
the \emph{walls}, in the space of pairs of marginals; the \emph{chambers} are the connected components of
the complement of the walls.
\end{definition}

\begin{theorem}[Coupling polytopes are Delzant]\label{thm:delzant}
If $(f,g)$ is nondegenerate, then $\Cpl(f,g)$ is a Delzant polytope with respect to $\Lambda$: it is
simple, its facet normals are rational, and at each vertex the primitive edge vectors form a basis of
$\Lambda$. Its edges at a vertex are the fundamental cycles of a spanning tree.
\end{theorem}
\begin{proof}
Regard $\alpha$ as a weighting of the edges of the complete bipartite graph on $Y\sqcup Y'$.
\emph{Vertices are forests.} If the support of a feasible $\alpha$ contains a cycle, adding $\pm\epsilon$
alternately around it stays feasible, so $\alpha$ is not a vertex. If the support is a forest, the
marginal equations determine $\alpha$ on it uniquely (peel off leaves), so $\alpha$ is a vertex.
\emph{Nondegeneracy makes them spanning trees.} Every node has positive marginal and hence an incident
edge of the support. If the support forest had two or more components, take one with row set $I$ and
column set $J$. Then $\sum_If=\sum_Jg$; $I$ and $J$ are nonempty, and if $I=Y$ then $\sum_Jg=1$ forces
$J=Y'$ and the component is everything. So $I$ and $J$ are proper, contradicting nondegeneracy.
\emph{Simplicity and smoothness.} At a vertex with spanning tree $T$, the zero cells are the
$nm-(n+m-1)=(n-1)(m-1)$ non-tree cells $c$. Each closes a unique cycle in $T\cup\{c\}$, giving a
circulation $C_c$ with entries in $\{0,\pm1\}$ and $C_c(c)=1$. A circulation $w$ equals
$\sum_{c\notin T}w(c)C_c$, because the difference is a circulation supported on the tree, hence zero.
So the feasible directions at the vertex form the cone $\{\sum_ct_cC_c:t\ge0\}$ on exactly
$\dim\Cpl(f,g)$ independent generators (simplicity), and the $C_c$ are integral and form a $\mathbb{Z}$-basis
of $\Lambda$, since integer coefficients $w(c)$ suffice (smoothness). The facet normals are the coordinate
functionals $\alpha\mapsto\alpha(c)$, which are rational.
\end{proof}

By Delzant's theorem \cite{Delzant1988} there is a compact connected toric symplectic manifold
$(X_{f,g},\omega)$ of real dimension $2(n-1)(m-1)$, unique up to equivariant symplectomorphism, whose
moment map $\mu\colon X_{f,g}\to\Cpl(f,g)$ has image the coupling polytope. The fibre over an interior
coupling is a Lagrangian torus $T^{(n-1)(m-1)}$, and the fibre over a vertex is a point.

\begin{proposition}[Fisher metric and Guillemin's potential]\label{prop:guillemin}
Suppose every inequality $\alpha(c)\ge0$ defines a facet of $\Cpl(f,g)$. The canonical K\"ahler metric
of $X_{f,g}$ constructed by Guillemin \cite{Guillemin1994} has, in action--angle coordinates over the
interior, the form $\mathrm{Hess}\,G$ in the polytope directions, with symplectic potential
\[
G(\alpha)=\tfrac12\sum_{c}\alpha(c)\log\alpha(c).
\]
Hence the Fisher metric on $\Cpl(f,g)$ is twice the restriction of the canonical K\"ahler metric to the
real section, and the product coupling $f\otimes g$, the coupling of maximal entropy, is the unique
critical point of $G$ on the polytope.
\end{proposition}
\begin{proof}
Guillemin's potential is $\tfrac12\sum_i\ell_i\log\ell_i$ over the facets, where $\ell_i\ge0$ are the
facet inequalities written with primitive inward normals. Here $\ell_c(\alpha)=\alpha(c)$, and the
coordinate functional is primitive on $\Lambda$ because some integer circulation has entry $1$ at $c$.
Then $\mathrm{Hess}\,G(v,v)=\tfrac12\sum_cv(c)^2/\alpha(c)$. The critical-point condition is that
$\log\alpha$ is orthogonal to all circulations, that is $\log\alpha(y,y')=r(y)+s(y')$, which together with
the marginals gives $\alpha=f\otimes g$. Uniqueness follows from strict convexity.
\end{proof}

If some cell is not a facet, for instance when $f(y)+g(y')>1$ forces $\alpha(y,y')>0$, the potential
omits that term and the identity holds only up to those terms.

\subsection{Walls, chambers and monotone cases}

\begin{proposition}[Chambers]\label{prop:chambers}
The set of spanning trees that support vertices, and hence the combinatorial type of $\Cpl(f,g)$ and
the toric manifold $X_{f,g}$ up to equivariant diffeomorphism, is constant on each chamber. The type
can change across a wall. For $f=(\tfrac12,\tfrac12)$ and $g=(t,\tfrac23-t,\tfrac13)$ the polytope is a
hexagon for $t<\tfrac12$ and a quadrilateral for $t>\tfrac12$, both Delzant: two facets become redundant
across the wall $g_1=f_1$.
\end{proposition}
\begin{proof}
The tree solution on a spanning tree $T$ assigns to each tree edge the value $\sum_If-\sum_Jg$ (up to
sign) for the bipartition $(I,J)$ obtained by deleting that edge. So feasibility of $T$ depends only on
the signs of such differences, which are constant on chambers. The example is checked in
\texttt{check\_conditionals.py}.
\end{proof}

\begin{proposition}[Monotone cases]\label{prop:monotone}
Assume $(f,g)$ is nondegenerate and every cell defines a facet. Then $X_{f,g}$ is monotone if and only
if $f$ and $g$ are uniform, and uniform marginals are nondegenerate exactly when $\gcd(n,m)=1$. In
particular, in the square case $Y=Y'$ of coupling 2-cells no nondegenerate pair is monotone. On the
wall, the uniform $2\times2$ pair still gives a Delzant segment, $X=\mathbb{CP}^1$, with the product
coupling as its midpoint; the uniform $n\times n$ pair with $n\ge3$ gives a polytope that is not simple.
\end{proposition}
\begin{proof}
For a Delzant polytope written with primitive normals, the toric manifold is monotone exactly when
some interior point has equal lattice distance to all facets: the classes of the toric divisors span
$H^2$ subject to the linear relations given by the normals, $c_1$ is the sum of the divisor classes, and
$[\omega]$ is the sum weighted by the facet distances. Here the distances are the entries
$\alpha(c)$, so they are equal exactly when $\alpha$ is constant, which forces $f$ and $g$ to be
uniform; then $\alpha=f\otimes g$. Uniform marginals are nondegenerate iff $a/n\ne b/m$ for
$0<a<n$, $0<b<m$, iff $\gcd(n,m)=1$, which fails when $n=m$. The two wall cases are checked in
\texttt{check\_conditionals.py} (T3).
\end{proof}

\begin{example}[The degree-6 del Pezzo surface]\label{ex:dp6}
For $f$ uniform on two points and $g$ uniform on three, $\Cpl(f,g)$ is a hexagon. In the basis of
$\Lambda$ above, its facet normals are $\pm e_1$, $\pm e_2$ and $\pm(e_1+e_2)$. This is the fan of
$\mathbb{CP}^2$ blown up at three points, the del Pezzo surface of degree $6$. All six entries of the
product coupling equal $\tfrac16$, so it is the equidistant point. The case $3\times4$ (uniform) gives a
Delzant polytope of dimension $6$ with $96$ vertices and all $12$ cells facets, again monotone.
\end{example}

%=====================================================================
\section{The imported layer: Floer theory of the special fibre}\label{sec:imported}
%=====================================================================

We now quote results from symplectic topology. They are not proved or checked here.

\paragraph{Nondisplaceability.} For a monotone toric manifold, the torus fibre over the point at equal
distance from all facets (the \emph{special fibre}) is a stem in the sense of Entov and Polterovich,
hence nondisplaceable \cite{EntovPolterovich2006}. For toric Fano manifolds, Cho and Oh show that a
torus fibre with a flat $U(1)$ connection has nonvanishing Floer cohomology when it corresponds to a
critical point of the disc potential \cite{ChoOh2006}; see \cite{Cho2004} for the Clifford torus and
\cite{FOOO2010} for general compact toric manifolds.

\paragraph{Consequence for couplings.} In the monotone cases of \cref{prop:monotone}, the special fibre
lies over the product coupling $f\otimes g$. For \cref{ex:dp6}, the disc potential of the monotone
special fibre is, up to the factor $T^{1/6}$, the Laurent polynomial
\[
W(x,y)=x+y+\frac1x+\frac1y+xy+\frac1{xy},
\]
read off from the facet normals, and it has a critical point at $(x,y)=(1,1)$ (checked exactly), so the
special fibre with trivial holonomy has nonzero Floer cohomology by \cite{ChoOh2006}.

The maximal coupling, by contrast, has zero entries whenever $f\ne g$, so it lies on the boundary of
$\Cpl(f,g)$, where the torus fibre collapses (\texttt{check\_conditionals.py}, T6). The first draft's
conjecture that the critical points are the maximal couplings should therefore be replaced. The
geometrically distinguished coupling is the product coupling, which is also the maximum-entropy
coupling and the critical point of Guillemin's potential (\cref{prop:guillemin}).

%=====================================================================
\section{Stability conditions: a question}\label{sec:stability}
%=====================================================================

This subsection replaces a draft that proposed a Bridgeland stability condition whose central charge
is the disc potential, whose heart is the committed states, and whose Harder--Narasimhan filtration
is the obstruction tower. None of those identifications is well typed. A central charge is an additive
map $K_0\to\mathbb{C}$ and $W$ is a function on a torus. Committed states are vertices of a polytope
and form no abelian category. The obstruction tower assigns numbers to a claim; it does not
filter an object. What can be said precisely is the following.

\paragraph{Bridgeland stability.} A stability condition on a triangulated category $\mathcal D$ is a
pair $(Z,\mathcal A)$. Here $\mathcal A$ is the heart of a bounded $t$-structure, and
$Z\colon K_0(\mathcal D)\to\mathbb{C}$ is additive, sends nonzero objects of $\mathcal A$ to the upper
half-plane together with $\mathbb{R}_{<0}$, and admits Harder--Narasimhan filtrations; the support
property is also required \cite{Bridgeland2007}. On any smooth projective surface $X$, and for
any ample class $\omega$, there is a stability condition with central charge
\[
Z_\omega(E)=-\int_Xe^{-i\omega}\operatorname{ch}(E)
\]
whose heart is a tilt of $\operatorname{Coh}X$ \cite{ArcaraBertram2013,MacriSchmidt2017}. (A $B$-field
can be added; we set it to zero.)

\begin{proposition}[Marginals give stability conditions on the $2\times3$ surfaces]\label{prop:stab}
Let $(f,g)$ be nondegenerate in the $2\times3$ family, let $X_\Sigma$ be the toric surface of its
chamber, and let $D_c$ be the toric divisor of the facet $\{\alpha_c=0\}$. Put
\[
[\omega_{f,g}]=\sum_{c\ \mathrm{facet}}f_ig_j\,[D_c],\qquad c=(i,j),
\]
the product coupling read as divisor coefficients. Then:
\begin{enumerate}[label=(\roman*)]
\item $[\omega_{f,g}]$ is the class of the toric symplectic form with moment polygon $\Cpl(f,g)$. It is
  independent of the base point used to write the polygon, and it is ample.
\item $\omega_{f,g}\cdot D_c$ is the lattice length of the edge $\{\alpha_c=0\}$ of the coupling polygon,
  and $\tfrac12\omega_{f,g}^2$ is its lattice area.
\item There is a stability condition $\sigma_{f,g}=(Z_{\omega_{f,g}},\mathcal A_{\omega_{f,g}})$ on
  $D^b(\operatorname{Coh}X_\Sigma)$. It depends continuously on $(f,g)$ within the chamber, and
  \[
  Z(\mathcal O_X)=\operatorname{area}\Cpl(f,g),\qquad Z(\mathcal O_x)=-1,\qquad
  Z(\mathcal O_{D_c})=\tfrac12D_c^2+i\,\operatorname{length}\{\alpha_c=0\}.
  \]
\item As $(f,g)$ approaches a wall from inside a chamber, one edge length tends to $0$, its divisor
  $E$ is a $(-1)$-curve, and $Z(\mathcal O_E)\to-\tfrac12$, a point on the negative real axis.
\end{enumerate}
\end{proposition}
\begin{proof}
(i) If a Delzant polygon is $\{u:\langle u,v_c\rangle+\lambda_c\ge0\}$, the symplectic class is
$\sum_c\lambda_c[D_c]$, up to the normalization of $\omega$. Here the facet functions are
$\alpha_c(u)=\langle u-u_0,v_c\rangle+\alpha_c(u_0)$ for any interior base point $u_0$, and choosing
$u_0=f\otimes g$ gives the formula. Another choice of base point changes the coefficients by
$\langle w,v_c\rangle$, which is a linear relation among the $[D_c]$. A Kähler class on a projective
toric surface is ample. (ii) These are standard toric identities: the degree of $\omega$ on $D_c$ is
the lattice length of the corresponding edge, and the symplectic volume is the lattice area
(Duistermaat--Heckman). (iii) Apply \cite{ArcaraBertram2013,MacriSchmidt2017} with $\omega=\omega_{f,g}$,
which varies continuously with $(f,g)$. The values follow from $\operatorname{ch}(\mathcal O_X)=1$,
$\operatorname{ch}(\mathcal O_x)=[\mathrm{pt}]$ and $\operatorname{ch}(\mathcal O_D)=D-\tfrac12D^2$. (iv) By
\cref{prop:chambers} and \texttt{check\_stability\_layer.py} (S5), the vanishing edge is that of a
$(-1)$-curve. All identities are checked in all $18$ chambers.
\end{proof}

So the probabilistic data enter Bridgeland stability through the Kähler class, and nowhere else. The
product coupling fixes the class, edge lengths of the coupling polygon are imaginary parts of central
charges, and its area is $Z(\mathcal O_X)$. At the uniform point $\omega=-\tfrac16K_X$ on the del Pezzo
surface, with area $\tfrac1{12}$ and six edges of length $\tfrac16$.

\paragraph{Questions.}
\begin{enumerate}[label=(Q\arabic*),leftmargin=*]
\item \emph{Across walls.} Crossing a wall blows up or blows down a point
  (\cref{prop:chambers}), and Orlov's theorem gives
  $D^b(\mathrm{Bl}_pX)=\langle\mathcal O_E(-1),D^b(X)\rangle$ \cite{Orlov1992}. As the wall is approached,
  $Z(\mathcal O_E)$ tends to the negative real axis. Does the path $\sigma_{f,g}$ extend across the wall to
  a path through stability conditions glued along Orlov's decomposition? If so, the chamber structure of
  the marginals becomes a wall structure in a single space of stability conditions.
\item \emph{Objects for couplings.} Is there a natural object $E_\alpha\in D^b(\operatorname{Coh}X)$ attached to
  an interior coupling $\alpha$? A candidate is the SYZ or mirror transform of the torus fibre over
  $\alpha$ with a flat connection. Does $Z(E_\alpha)$ have a probabilistic meaning, for instance in terms of
  the Fisher metric of \cref{prop:guillemin}?
\item \emph{Quantum corrections.} Iritani's integral structure identifies quantum-corrected central
  charges of objects of $D^b(\operatorname{Coh}X)$ with oscillatory integrals $\int e^{-W/z}$ of the
  mirror potential over integral cycles, and their asymptotics involve the $\widehat\Gamma$-class
  \cite{Iritani2009,GGI2016}. This is the precise sense in which the disc potential is related to
  central charges: through oscillatory integrals, not through its values. Do the corrected central
  charges, with $\omega=\omega_{f,g}$, define stability conditions on these surfaces? Does the monotone
  point $f\otimes g$ uniform play a special role?
\item \emph{Higher dimensions.} For $(n-1)(m-1)\ge3$ the toric manifolds have complex dimension at least
  three. There, the existence of stability conditions is known only in special cases, through
  Bogomolov--Gieseker type inequalities \cite{BMT2014}. The construction of \cref{prop:stab} therefore
  extends only conditionally.
\end{enumerate}

\paragraph{What is not claimed.} We make no identification of the disc potential with a central
charge, of committed states with a heart, of the obstruction tower with Harder--Narasimhan filtrations,
or of hallucination with instability. The obstruction tower and the stability conditions above live
on different objects (claims and admissibility data on one side, coherent sheaves on $X_{f,g}$ on the
other), and any relation between them would have to be constructed.

%=====================================================================
\section{The conjectural layer}\label{sec:conjectural}
%=====================================================================

\begin{proposition}[The $2\times3$ family {\cite[\S3]{Awasthi-Family}}]\label{prop:family23}
In the $2\times3$ family there are $18$ chambers. Every toric surface $X_{f,g}$ is the degree-6 del Pezzo surface or a
toric blow-down of it, and crossing one wall is exactly one toric blow-up or blow-down.
\end{proposition}

\begin{conjecture}[Family over the chamber complex]\label{conj:family}
For general $n\times m$, the toric manifolds $X_{f,g}$ form a family over each chamber, and crossing a wall is a toric
birational transformation (in higher dimension possibly a flip) determined by the pair $(I,J)$ defining the wall. On a
wall, $X_{f,g}$ degenerates to a singular toric variety whose moment polytope is the non-simple $\Cpl(f,g)$.
\end{conjecture}

\begin{question}[Location of Floer-nontrivial fibres]\label{q:location}
Outside the monotone cases, the fibres with nonvanishing Floer cohomology (possibly after bulk
deformation) are located by the potential of \cite{FOOO2010}. Is the product coupling always among
them, and do they have a statistical meaning?
\end{question}

\begin{question}[Beyond the minimal cone]\label{conj:ainfty}
\Cref{thm:sqzero} makes whiskering an $A_\infty$-functor into the square-zero cone, with two components. (i) In the free cone,
do the higher Boolean cumulants give the higher components? (ii) Does the construction extend to the
multi-object setting, an $A_\infty$-functor between linear categories whose objects are 1-cells, compatible with the
nerve of \cref{sec:nerve}?
\end{question}

\begin{conjecture}[Fukaya-type functor]\label{conj:fukaya}
There is a functor from a category built from the Fukaya categories of the $X_{f,g}$ (with special
fibres as distinguished objects) to a linearization of the coupling 2-cells, compatible with gluing on
one side and composition of Lagrangian correspondences on the other.
\end{conjecture}

%=====================================================================
\section{Next steps}\label{sec:next}
%=====================================================================

The companion \cite{Awasthi-Family} computes every layer explicitly: the $2\times3$ family, the square $2\times2$ case,
levels two and three, disc potentials, stability, cumulants, the nerve, the tower fibration and the family over
the marginals. The open problems are:
\begin{enumerate}[label=(\arabic*),leftmargin=*]
\item \emph{Across a wall.} Decide (Q1) of \cref{sec:stability} for the wall $f_1=g_1$ between the pentagon and
  hexagon chambers.
\item \emph{General wall-crossing.} Prove \cref{conj:family} beyond $2\times3$.
\item \emph{Higher horns.} Give explicit unfillable $n$-horns and their contextual fractions for all $n\ge5$, and
  compute the best approximate fillers.
\item \emph{Multi-object cone.} Answer \cref{conj:ainfty}(ii), linking the cone to the nerve.
\item \emph{Fukaya.} \Cref{conj:fukaya}.
\end{enumerate}

\section{Computational checks}\label{sec:code}

All checks pass.
\begin{itemize}[leftmargin=*]
\item \texttt{check\_glue.py} (30 checks): \cref{prop:1cells,thm:witness}.
\item \texttt{check\_stochastic\_basechange.py} (13 checks): \cref{thm:fibre}; the hexagon of \cref{ex:dp6}; the Delzant
  property for random nondegenerate $3\times3$ marginals; non-simplicity of the uniform $3\times3$ case; the Hessian
  identity and critical point of \cref{prop:guillemin}.
\item \texttt{check\_conditionals.py} (7 checks): \cref{thm:delzant,prop:chambers,prop:monotone,ex:dp6}.
\item \texttt{check\_stability\_layer.py} (6 checks): \cref{prop:stab} in all $18$ chambers of the $2\times3$ family.
\item In the companion \cite{Awasthi-Family}:
  \begin{itemize}[leftmargin=*]
  \item \texttt{check\_tower\_fibration.py} (7 checks): \cref{thm:fibre,prop:gradefibre};
  \item \texttt{check\_nerve.py} (10 checks): \cref{thm:2horn,thm:higherhorn,prop:boundaryfibre,thm:sqzero};
  \item \texttt{check\_cumulants.py} (8 checks): \cref{prop:kappa2};
  \item \texttt{check\_polytope\_companion.py}, \texttt{check\_family.py} and \texttt{check\_family\_geometry.py}:
    \cref{prop:family23}.
  \end{itemize}
\end{itemize}

\appendix
%=====================================================================
\section{Corrections to the first draft}\label{sec:corrections}
%=====================================================================

\begin{enumerate}[leftmargin=*]
\item \emph{Candidate disc potential.} $W_\beta(f)=\sum_g\int_{\Cpl(f,g)}e^{-\beta c(\alpha)}\,d\nu$ sums over a
  continuum of 1-cells $g$ and is not defined as written. Its Laplace limit
  $\min_{g,\alpha}c(\alpha)$ is always $0$, attained at $g=f$, $\alpha=\Delta_f$, so it carries no
  information. The toric disc potential of $X_{f,g}$ is the correct object (\cref{sec:imported}).
\item \emph{Symplectic structure.} A symplectic form is antisymmetric and the Fisher metric is
  symmetric, so the restriction of the one cannot be the other. The correct statement is
  \cref{prop:guillemin}: the Fisher metric is twice the canonical K\"ahler metric on the real section
  of $X_{f,g}$.
\item \emph{Critical points.} ``Critical points of $\pi\colon E\to\Kl(\D)$'' is undefined, since $\pi$ is a
  functor, not a smooth map. The replacement is the special fibre over the product coupling in the
  monotone cases, and the walls of \cref{def:walls} as the locus where the family degenerates.
\item \emph{Monodromy and Dehn twists.} The family $(f,g)\mapsto X_{f,g}$ is constant on chambers
  and changes by toric birational modifications across walls (\cref{prop:chambers}); no vanishing
  cycle or Dehn twist has been identified. \Cref{conj:family} restates the question.
\item \emph{The $m_3$ conjecture.} Gluing is strictly associative (\cref{thm:witness}(i)), so each
  hom-category, regarded as an $A_\infty$-category, has $m_k=0$ for $k\ge3$. The whiskering defect
  compares $u\cdot(\beta\cdot\alpha)$ with $(u\cdot\beta)\cdot(u\cdot\alpha)$. Its inputs are two 2-cells and a
  1-cell, and it measures the failure of whiskering to be a functor. In $A_\infty$ language that is the
  second component $\mathcal{F}_2$ of an $A_\infty$-functor, not a ternary composition
  (now \cref{thm:sqzero}).
\item \emph{Homotopy category.} The graded shadow has hom-sets $[d(f,g),\infty)$ under addition. These
  are not the degree-zero cohomology of cochain complexes, so the shadow is not a homotopy category of
  an $A_\infty$-category without further structure.
\item \emph{Fisher versus Fisher--Rao.} These are the same metric (\cref{prop:interior}).
\item \emph{References.} The companion tower is now version 9; its posted earlier version has the
  title ``\ldots and Weak Distributive Laws'', a claim withdrawn in the revisions.
\end{enumerate}

\begin{thebibliography}{99}

\bibitem{Awasthi-Attribution}
V.~K. Awasthi.
\newblock Precise attribution of model interventions: attribution sets, mechanism reach, path defects
and gluing.
\newblock Technical report, Johns Hopkins University, October 2026. Draft 4.

\bibitem{Awasthi-Tower}
V.~K. Awasthi.
\newblock The three-level posterior tower: hierarchical {B}ayesian inference and mixed distributive laws.
\newblock Technical report, Johns Hopkins University, October 2026. Revised version 9. An earlier
version is posted as DOI 10.13140/RG.2.2.29743.50088.

\bibitem{ArcaraBertram2013}
D.~Arcara and A.~Bertram.
\newblock Bridgeland-stable moduli spaces for {$K$}-trivial surfaces (with an appendix by M.~Lieblich).
\newblock \emph{Journal of the European Mathematical Society}, 15(1):1--38, 2013.

\bibitem{Awasthi-Family}
V.~K. Awasthi.
\newblock Coupling polytopes and toric geometry: the $2\times3$ family, the square case, higher levels and stability.
\newblock Companion paper, Johns Hopkins University, October 2026.

\bibitem{BMT2014}
A.~Bayer, E.~Macr\`i, and Y.~Toda.
\newblock Bridgeland stability conditions on threefolds {I}: {B}ogomolov--{G}ieseker type inequalities.
\newblock \emph{Journal of Algebraic Geometry}, 23(1):117--163, 2014.

\bibitem{Bridgeland2007}
T.~Bridgeland.
\newblock Stability conditions on triangulated categories.
\newblock \emph{Annals of Mathematics}, 166(2):317--345, 2007.

\bibitem{Cho2004}
C.-H. Cho.
\newblock Holomorphic discs, spin structures, and {F}loer cohomology of the {C}lifford torus.
\newblock \emph{International Mathematics Research Notices}, 2004(35):1803--1843, 2004.

\bibitem{ChoOh2006}
C.-H. Cho and Y.-G. Oh.
\newblock Floer cohomology and disc instantons of {L}agrangian torus fibers in {F}ano toric manifolds.
\newblock \emph{Asian Journal of Mathematics}, 10(4):773--814, 2006.

\bibitem{Delzant1988}
T.~Delzant.
\newblock Hamiltoniens p\'eriodiques et images convexes de l'application moment.
\newblock \emph{Bulletin de la Soci\'et\'e Math\'ematique de France}, 116(3):315--339, 1988.

\bibitem{EntovPolterovich2006}
M.~Entov and L.~Polterovich.
\newblock Quasi-states and symplectic intersections.
\newblock \emph{Commentarii Mathematici Helvetici}, 81(1):75--99, 2006.

\bibitem{FOOO2010}
K.~Fukaya, Y.-G. Oh, H.~Ohta, and K.~Ono.
\newblock Lagrangian {F}loer theory on compact toric manifolds {I}.
\newblock \emph{Duke Mathematical Journal}, 151(1):23--174, 2010.

\bibitem{GGI2016}
S.~Galkin, V.~Golyshev, and H.~Iritani.
\newblock Gamma conjectures for {F}ano manifolds.
\newblock \emph{Duke Mathematical Journal}, 165(11):2005--2077, 2016.

\bibitem{Guillemin1994}
V.~Guillemin.
\newblock Kaehler structures on toric varieties.
\newblock \emph{Journal of Differential Geometry}, 40(2):285--309, 1994.

\bibitem{Iritani2009}
H.~Iritani.
\newblock An integral structure in quantum cohomology and mirror symmetry for toric orbifolds.
\newblock \emph{Advances in Mathematics}, 222(3):1016--1079, 2009.

\bibitem{MacriSchmidt2017}
E.~Macr\`i and B.~Schmidt.
\newblock Lectures on {B}ridgeland stability.
\newblock In \emph{Moduli of Curves}, Lecture Notes of the Unione Matematica Italiana 21. Springer,
2017. arXiv:1607.01262.

\bibitem{Orlov1992}
D.~O. Orlov.
\newblock Projective bundles, monoidal transformations, and derived categories of coherent sheaves.
\newblock \emph{Izvestiya Rossiiskoi Akademii Nauk. Seriya Matematicheskaya}, 56(4):852--862, 1992.

\bibitem{AbramskyBrandenburger2011}
S.~Abramsky and A.~Brandenburger.
\newblock The sheaf-theoretic structure of non-locality and contextuality.
\newblock \emph{New Journal of Physics}, 13:113036, 2011.

\bibitem{BFMY1983}
C.~Beeri, R.~Fagin, D.~Maier, and M.~Yannakakis.
\newblock On the desirability of acyclic database schemes.
\newblock \emph{Journal of the ACM}, 30(3):479--513, 1983.

\bibitem{DeLoeraOnn2004}
J.~A. De~Loera and S.~Onn.
\newblock All rational polytopes are transportation polytopes and all polytopal integer sets are contingency
tables.
\newblock In \emph{Integer Programming and Combinatorial Optimization (IPCO 2004)}, Lecture Notes in Computer
Science 3064. Springer, 2004.

\bibitem{HPT1}
G.~C. Drummond-Cole, J.-S. Park, and J.~Terilla.
\newblock Homotopy probability theory {I}.
\newblock \emph{Journal of Homotopy and Related Structures}, 10(3):425--435, 2015.

\bibitem{HPT2}
G.~C. Drummond-Cole, J.-S. Park, and J.~Terilla.
\newblock Homotopy probability theory {II}.
\newblock \emph{Journal of Homotopy and Related Structures}, 10(3):623--635, 2015.

\bibitem{LermanTolman1997}
E.~Lerman and S.~Tolman.
\newblock Hamiltonian torus actions on symplectic orbifolds and toric varieties.
\newblock \emph{Transactions of the American Mathematical Society}, 349(10):4201--4230, 1997.

\bibitem{PachterSturmfels2005}
L.~Pachter and B.~Sturmfels, editors.
\newblock \emph{Algebraic Statistics for Computational Biology}.
\newblock Cambridge University Press, 2005.

\bibitem{SpeicherWoroudi1997}
R.~Speicher and R.~Woroudi.
\newblock Boolean convolution.
\newblock In \emph{Free Probability Theory}, Fields Institute Communications 12. American Mathematical
Society, 1997.

\bibitem{Vorobev1962}
N.~N. Vorob'ev.
\newblock Consistent families of measures and their extensions.
\newblock \emph{Theory of Probability and Its Applications}, 7(2):147--163, 1962.

\end{thebibliography}
\end{document}
